\chapter{Introduction}
\label{chap:intro}

Data streams arrive continuously and without a known end, and the process that generates them can change while they are observed. Clustering such a stream requires an algorithm that processes each sample as it arrives, keeps a bounded summary of what it has seen and follows the changes of the data, a family of methods known as stream or evolving clustering \cite{aggarwal2003framework,cao2006density,wang2023data}. \MicroTEDAClus \cite{maia2020evolvclus} is one such algorithm. It summarizes the stream into microclusters by the eccentricity and the typicality of the TEDA framework \cite{angelov2014outside}, which assume no distribution for the data and are computed recursively, and it groups the microclusters into macroclusters that represent the clusters of the stream. Its authors report that it performs on par with established stream clustering algorithms, and that it adapts better than them to concept drift.

This work reassesses \MicroTEDAClus and proposes an improved model. The reassessment finds issues in the formulation of the algorithm, differences between the formulation and the implementation that produced the published results, and flaws in the methodology of the original experiments, among them preprocessing steps that leak information about the stream and an evaluation that predicts each batch after learning it. Under a methodology that corrects these flaws, \MicroTEDAClus is frequently outperformed by the benchmark models. The improved model, \MicroTEDAClusPlus, is the result of controlled experiments on corrections and modifications of the algorithm. On the synthetic datasets of the original work it improves on \MicroTEDAClus in clustering quality, execution time and memory, and it is among the best benchmark models in the datasets with concept drift. In an experiment modernized in its environment, its benchmark models and its selection of hyperparameters, it remains indistinguishable from the best model in most datasets. On real datasets, however, it is among the worst models. A second effort of controlled experiments on those datasets produces the final version of \MicroTEDAClusPlus, which is no longer among the worst models in most real datasets and is the best in one of them, while it remains behind the best models on the real streams with concept drift.

\section{Motivation}
\label{sec:motivation}

The results of \MicroTEDAClus are the basis of later work. The doctoral thesis of \cite{galindo2021lastmile} updates the algorithm and applies it to the last-mile problem in logistics, assuming the validity of the original evaluation without examining it. A model whose evaluation supports further work must be verified, and the verification of empirical results in machine learning is not a formality: information leaked from the data into the evaluation, among other errors, has been found to inflate the published results of many studies \cite{kaufman2012leakage,kapoor2023leakage}. The algorithm is also attractive for what it promises. The TEDA metrics require neither a distribution for the data nor a predefined number of clusters, and an algorithm built on them that is accurate, fast and robust to concept drift would be a useful tool for stream clustering.

The datasets of the original work are synthetic streams and benchmarks of static shapes, whose clusters, drifts and noise are known by construction. A model that is to serve applications must also be judged on data recorded from physical processes, whose classes are not drawn from chosen distributions and whose drifts come from the conditions of the recording. An evaluation on real data tests whether the conclusions reached on synthetic data hold.

This work verifies the claims of \MicroTEDAClus before building on them, and improves the model where the verification shows weaknesses. The TEDA framework is that of \cite{angelov2014outside}. The formulation of \MicroTEDAClus and its reference implementation are those of \cite{maia2020evolvclus}, and its updated formulation is that of \cite{galindo2021lastmile}. The benchmark models are those of their respective authors. In the main experiment they are executed through the implementations of the \codename{River} library \cite{montiel2021river} or, for \CEDAS, \TEDACloud and \AutoCloud, through implementations made for this work from the official code or from the descriptions in their papers; in the modernized and real-data experiments, through \codename{River}, \codename{scikit-learn} \cite{pedregosa2011scikit} and the \codename{bico} package \cite{fichtenberger2013bico}. The real datasets are those of the sources cited where they are introduced. The reassessment of the formulation, of the implementation and of the methodology, the methodology that replaces the original one, the vectorized implementation, the experiments of the improvement and \MicroTEDAClusPlus, the modernized experiment with its ported stream generators and its search ranges by rule, the real-data experiment, and the second effort of improvement with the final version of \MicroTEDAClusPlus are the contribution of this work.

\section{Objectives}
\label{sec:objectives}

From the motivation presented in the previous section, the objectives of this work are:

\begin{enumerate}[a)]
\item to examine the formulation of \MicroTEDAClus and its public implementation, and identify their issues and the differences between them;
\item to examine the methodology of the experiments of the original work, and define a methodology that evaluates every model alike and without its flaws;
\item to evaluate \MicroTEDAClus against benchmark models under that methodology;
\item to propose an improved model by controlled experiments on the issues identified, and evaluate it under the same methodology against \MicroTEDAClus and the benchmark models;
\item to modernize the experiment, with a single environment, benchmark models from maintained libraries, a selection of hyperparameters that depends neither on the scale of the data nor on the speed of the machine, and the cost of every model measured with each task executed alone, and to evaluate the improved model under it;
\item to evaluate the improved model on real datasets and, where it falls short, improve it again by controlled experiments, and compare the result with the benchmark models.
\end{enumerate}

\section{Structure of the Thesis}
\label{sec:structure}

The work is divided into nine chapters besides this introduction. \Cref{chap:background_relatedwork} presents the TEDA framework, the formulation of \MicroTEDAClus and its update, together with the issues of both. \Cref{chap:methodology_reassessment} examines the implementation of \MicroTEDAClus and the methodology of the original experiments, defines the methodology of this work and evaluates \MicroTEDAClus and the benchmark models under it. \Cref{chap:ModelImprovement} optimizes the implementation, evaluates corrections and improvements of the algorithm, defines \MicroTEDAClusPlus and compares it with \MicroTEDAClus in the main experiment. \Cref{chap:results} places \MicroTEDAClusPlus among the benchmark models. \Cref{chap:modernization} replaces the environment, the benchmark models and the selection of hyperparameters, and repeats the main experiment under them. \Cref{chap:real_data} evaluates the models on real datasets. \Cref{chap:real_improvement} diagnoses \MicroTEDAClusPlus on those datasets, improves it again and compares its final version with the benchmark models. \Cref{chap:discussion} reads the results of these chapters together, and \cref{chap:conclusions} summarizes the findings, their limits and the directions of future work. \Cref{app:mi_tables} collects complementary tables of the experiments of \cref{chap:ModelImprovement}.

\clearpage
